\textit{Generalización temporal}. Para verificar que el desempeño no depende de una partición particular dentro de 2019, el marco se entrenó con los inviernos previos y se aplicó a un invierno completo no visto (2018, 2019; Tabla~\ref{tab_crosswinter}). Con $h=1$, el mejor modelo superó al Naive en 6 de 8 combinaciones estación--invierno, con un MAE relativo entre 0.81 y 1.09. \textit{Generalización espacial}. Al entrenar el componente estocástico en una estación y aplicarlo a otra (Tabla~\ref{tab_crossstation}, $h=1$), el MAE fue en promedio 1.05 veces el del modelo local (rango 0.80--1.45); el modelo transferido superó al Naive de la estación destino en 12 de 36 casos. El modelo entrenado con las cuatro estaciones combinadas obtuvo un MAE relativo medio de 0.94 frente al local. Esto indica en qué medida la dinámica de corto plazo del residuo es compartida por la red y, por tanto, transferible a estaciones sin historial suficiente.
